\section{国际税收法律制度}
\begin{enumerate}
  \item 3 对概念和 1 个制度。考点很固定。
  \item 税收管辖权。
        \begin{enumerate}
          \item 居民税收管辖权。属人管辖向税收领域的延伸。境内外所得。\zd{无限纳税义务}。
          \item 来源地税收管辖权。属地管辖权。一般是对\zd{非居民}。4 种。\zd{营业、投资、劳务和财产所得}。
        \end{enumerate}
  \item 各国基本上适用“常设机构”（\zd{分支机构、办事处、工厂、矿场等}）原则对营业所得从源征税。
  \item 国际重复征税 vs 国际重叠征税。先后被征。\zd{1 个倒霉蛋叫重复}，2 个倒霉蛋叫重叠。
  \item 国际避税 vs 国际逃税。合法手段叫避税，非法手段叫逃税。
  \item 双边税收协定。
  \item CRS/金融账户涉税信息\zd{自动}交换标准。共同申报准则。越来越重要的考点。
        \begin{enumerate}
          \item 各国承担义务的税务机关\zd{自动交换，1 年 1 次}。
          \item 以\zd{税收居民}（而非国籍）作为识别依据。
          \item 交换信息仅限\zd{金融账户信息}。
          \item 需要申报并交换的只有\zd{产生现金流}的资产和有\zd{现金价值}的金融资产。\al{存款账户、托管账户、股权、基金、信托、具有现金价值的保险合同以及年金合同}。
          \item 投资海外\zd{房产、珠宝、艺术品、贵金属}等无需申报。
        \end{enumerate}
\end{enumerate}
